\vspace{1em}
% ------------------------------------------------------------------------------
% Conclusão
% ------------------------------------------------------------------------------


\chapter{Conclusão}
\label{chap_conclusao}

Os resultados obtidos nesta prática mostram que o random walk simples em 1D é um modelo adequado para estudar, de forma objetiva, a relação entre simulação estocástica e probabilidade teórica.

A simulação de 10 caminhadas com 100 passos permitiu observar trajetórias distintas e deslocamentos finais variados, como esperado em processos aleatórios. A análise teórica por meio de
\begin{equation}
P_n(d)=\frac{1}{2^n}\binom{n}{\frac{n+d}{2}}
\end{equation}
permitiu atribuir probabilidade aos deslocamentos observados no experimento.

Com o reforço de Monte Carlo (20000 simulações), verificou-se boa concordância entre distribuição empírica e teórica, evidenciando que a implementação computacional está consistente com o modelo matemático.

Conclui-se, portanto, que a prática atingiu seus objetivos de:
\begin{itemize}
    \item simular trajetórias de passeio aleatório;
    \item analisar deslocamentos finais;
    \item validar resultados numéricos com a formulação probabilística analítica.
\end{itemize}

Como continuidade, o mesmo procedimento pode ser estendido para estudar random walks com viés ($p\neq 1/2$), maior número de dimensões e estatísticas de primeira passagem.

